
\documentclass[11pt,a4paper]{article}

\usepackage[T1]{fontenc}
\usepackage[utf8]{inputenc}
\usepackage[french]{babel}

\usepackage{amsmath}
\usepackage{amssymb}
\usepackage{amsthm}
\usepackage{mathtools}
\usepackage{hyperref}
\usepackage{cleveref}
\usepackage{microtype}

\newtheorem{theorem}{Théorème}[section]
\newtheorem{lemma}[theorem]{Lemme}
\newtheorem{proposition}[theorem]{Proposition}
\newtheorem{corollary}[theorem]{Corollaire}

\theoremstyle{definition}
\newtheorem{definition}[theorem]{Définition}

\theoremstyle{remark}
\newtheorem{remark}[theorem]{Remarque}

\title{Research Notes on the Riemann Hypothesis}
\author{Jason ALOYAU}
\date{\today}

\begin{document}

\maketitle

\begin{abstract}
Ce document constitue un cadre de recherche,
de vérification et de documentation relatif à
l'hypothèse de Riemann.

Aucune démonstration de l'hypothèse de Riemann
n'est revendiquée dans ce manuscrit.
\end{abstract}

\tableofcontents


\section{Statement of the Problem}

\begin{definition}
For $\Re(s)>1$,

\[
\zeta(s)
=
\sum_{n=1}^{\infty}\frac{1}{n^s}.
\]

\end{definition}

The function admits an analytic continuation to the entire complex plane except for a simple pole at $s=1$.

\begin{theorem}[Riemann Hypothesis]
Every non-trivial zero $\rho$ of $\zeta(s)$ satisfies

\[
\Re(\rho)=\frac12.
\]

\end{theorem}

The vertical line

\[
\Re(s)=\frac12
\]

is called the critical line.

The domain

\[
0<\Re(s)<1
\]

is called the critical strip.


\section{Known Results}

The following results are established.

\subsection{Euler Product}

For $\Re(s)>1$,

\[
\zeta(s)
=
\prod_{p}
\left(
1-p^{-s}
\right)^{-1}.
\]

This identity links the zeta function to prime numbers.

\subsection{Functional Equation}

The zeta function satisfies

\[
\zeta(s)
=
2^s
\pi^{s-1}
\sin\left(
\frac{\pi s}{2}
\right)
\Gamma(1-s)
\zeta(1-s).
\]

\subsection{Prime Number Theorem}

The prime-counting function satisfies

\[
\pi(x)
\sim
\frac{x}{\log x}.
\]

\subsection{Hardy Theorem}

There exist infinitely many zeros on the critical line.

\subsection{Computational Evidence}

The first very large collections of zeros that have been checked numerically lie on the critical line.

These verifications remain finite and therefore do not constitute a proof.

\section{Candidate Research Strategies}

This section collects possible directions of investigation.

\subsection{Equivalent Criteria}

Several statements are known to be equivalent to RH.

Examples include:

\begin{itemize}
\item Li Criterion.
\item Weil Criterion.
\item Nyman-Beurling Criterion.
\item Báez-Duarte Criterion.
\end{itemize}

A possible strategy is to simplify one of these formulations.

\subsection{Spectral Approaches}

The Hilbert--Pólya philosophy suggests that zeros could correspond to eigenvalues of a self-adjoint operator.

Finding such an operator remains open.

\subsection{Positivity Approaches}

Many equivalent criteria can be reformulated as positivity conditions.

The objective is to identify structures admitting rigorous verification.

\subsection{Formal Verification}

Incremental formal verification may reduce logical errors and hidden assumptions.
\section{Lemma Registry}

This section contains only verified intermediate results.

\begin{lemma}[Template]
Let $A$ be a mathematical object satisfying hypotheses H.

Then property P holds.
\end{lemma}

\begin{proof}
Proof placeholder.

The proof must explicitly justify all
limits, summations, integrations, and
analytic continuations.

No proof is provided at this stage.
\end{proof}

\subsection{Open Lemmas}

The following identifiers are reserved:

\begin{itemize}
\item LEM-0001
\item LEM-0002
\item LEM-0003
\end{itemize}

Their statements will be added as the project develops.
\section{Failure Analysis}

Many alleged proofs of RH fail for recurring reasons.

\subsection{Circular Reasoning}

An argument may unknowingly assume a result equivalent to RH.

\subsection{Invalid Exchanges}

Common issues involve unjustified exchanges between:

\[
\lim
\quad\text{and}\quad
\sum
\]

or

\[
\int
\quad\text{and}\quad
\sum.
\]

\subsection{Uniformity Problems}

An estimate valid for fixed parameters may fail under a limit process.

\subsection{Analytic Continuation Issues}

Extending a function beyond its domain of convergence requires full justification.

\subsection{Numerical Overreach}

Finite numerical verification is not equivalent to a proof of an infinite statement.

\section{Computational Evidence}

Computations play an important role in research.

\subsection{Objectives}

\begin{itemize}
\item Validate software implementations.
\item Explore numerical behaviour.
\item Detect counterexamples.
\item Generate conjectures.
\end{itemize}

\subsection{Limitations}

No finite computation can prove RH.

Computations only provide information about finitely many cases.

\subsection{Experiment Registry}

Current experiments include:

\begin{itemize}
\item EXP-001 - Zeta Sanity Checks.
\end{itemize}

Further experiments may investigate:

\begin{itemize}
\item Gram points.
\item Hardy Z-function.
\item Li coefficients.
\item Jensen polynomials.
\item Random matrix theory.
\end{itemize}


\section{Conclusion}

The Riemann Hypothesis remains one of the central unsolved problems of mathematics.

This repository is intended as:

\begin{itemize}
\item a research framework;
\item a reproducible computational environment;
\item a documentation platform;
\item a proof-management system;
\item a preparation ground for formal verification.
\end{itemize}

Future work will focus on:

\begin{enumerate}
\item equivalent formulations;
\item spectral approaches;
\item formal verification;
\item reproducible numerical experiments.
\end{enumerate}

At the current stage no proof of RH is claimed.


\bibliographystyle{alpha}
\bibliography{references}

\end{document}
